\section{Introduction}\label{section:intro}
Raft is a consensus algorithm designed for understandability and ease of implementation compared to the Paxos algorithm, and thus, it is widely utilized in various products to resolve consensus issues. A Raft cluster requires $2f+1$ servers to tolerate $f$ faults while remaining writable. This typically requires an odd number of replicas, at least 3, in many Raft-based systems. 

Most research efforts have focused on enhancing Raft's performance and scalability (BW-Raft, FlexiRaft, etc.) which are beneficial for large-scale clusters. However, there remains interest in reducing the number of servers without compromising fault tolerance. This is particularly important for organizations in retail and manufacturing industries which often need to manage thousands of small clusters. Reducing even a single server for each cluster can lead to significant cost savings. Diego Ongaro discussed replacing some servers with auxiliary ones (called witness servers) that store log entries for any failed server until recovery or replacement. Jehan-François Pâris and Darrell D. E. Long extend the idea to combine dynamic quorum and witness servers to achieve similar availability with fewer servers. 

To reduce costs, witness servers usually run in low configuration hardware, this can be performance bottleneck because witness needs to participate in the log replication when forming quorum. Additionally, heterogeneous server configurations increase deployment and maintenance complexity. Witness servers can also be implemented as secondary processes to leverage the power of regular servers. But this requires external mechanisms to coordinate their accessing to the witness data saved in shared storage, adding unnecessary external dependency and implementation complexity that limits its usage in real-world applications.

We hereby introduce a variant of Raft algorithm, called \textit{Extended Raft}, that allows using a shared storageo instead of a server as witness. \textit{Extended Raft} offers the following advantages that addressed the limitations of existing solutions:
\begin{enumerate}
    \item Cluster performance is independent of witness storage performance. Witness storage does not participate in the log replication in normal case, and only be accessed once for each failure occurrence or election. This guarantees cluster performance regardless of witness storage' hardware configuration or even availability in cerntain cases. This is especially useful for edge clusters that do not have stable network connection to external witness storage.
    \item Wide range of witness storage options. The witness storage can be a SMB/NFS share, a cloud storage, or even a USB stick plugged in the router, as long as it supports transactional read-modify-write across storage clients. One implementation for this requirement is to use optimistic concurrency control (OCC) in writing data. OCC is supported by most cloud storage services. We will show in this paper how to implement this using file share that supports hard link.
    \item More cost-effective. Shared storage is usually cheaper than a server, and can be shared among multiple clusters. This reduces the cost of witness storage and the overall cost of the system.
    \item Simplified deployment and maintenance. Managing shared storage is generally eaiser and more straightforward than managing a server because it involves fewer complex configurations, reduces the need for extensive technical expertise, and minimizes the ongoing maintenance tasks typically associated with server management.
\end{enumerate}
